
Il
résulte en particulier de 1.2.3 que pour tout Module gradué libre de type fini \(\underline{M}\), le nombre d'éléments d'une base homogène de \(\underline{M}\) ayant un degré donné ne dépend pas de la base choisie. En particulier, les entiers \(n_0\) et \(n_1\) intervenant dans la décomposition (1.2.2), définis par les conditions \(k_{n_0}\neq0\), \(k_{n_1}\neq0\), ne dépendent pas de la décomposition

\textbf{Proposition 1.3.} \uline{Soient} \(S\) \uline{un schéma quasi-compact}, \(\underline{E}\) \uline{un} \(\underline{O}_S\)\uline{-Module localement libre de type fini}, \(X\) \uline{le fibré projectif associé}. \uline{Soit} \(\underline{F}\) \uline{un} \(\underline{O}_X\)\uline{-Module de présentation finie}, \uline{plat sur} \(S\). \uline{Il existe un entier} \(n_0\), \uline{tel que pour} \(n\geq n_0\), \uline{on ait} :
\[
\begin{array}{ll}
\text{i)}& R^i f_*(\underline{F}(n))=0\quad\text{\uline{pour} }i\geq 1;\\[2mm]
\text{ii)}& f_*(\underline{F}(n))\text{ \uline{est localement libre de type fini sur} }S;\\[2mm]
\text{iii)}& \underline{M}=\displaystyle\bigoplus_{n\geq n_0} f_*(\underline{F}(n))\text{ \uline{est un} }\Sym(\underline{E})\text{\uline{-Module gradué de présentation finie} (1.2).}
\end{array}
\]

Les propriétés cherchées sont de nature locale sur \(S\); comme \(S\) est quasi-compact, on peut, en prenant le plus grand des entiers trouvés, se ramener au cas où \(S\) est affine.

Supposons d'abord \(S\) noethérien.

La propriété i) résulte alors du théorème de Serre.

La propriété ii) résulte de EGA III 7.9.10, \(n_0\) étant choisi tel que la propriété i) soit vérifiée.

D'après EGA II 3.4.5, il existe un \(\Sym(\underline{E})\)-Module gradué quasi-cohérent \(\underline{M}'\), de type fini, tel que \(\underline{F}=\Projz(\underline{M}')\). Par ailleurs, \(\Gamma_*(\underline{F})=\Gamma_*(\Projz \underline{M}')=\displaystyle\bigoplus_{n\in\mathbb Z} f_*(\underline{F}(n))\) est TN-isomorphe à \(\underline{M}'\) d'après EGA III 2.3.1; donc il existe un entier \(n_1\) tel que \(\displaystyle\bigoplus_{n\geq n_1} f_*(\underline{F}(n))\simeq \bigoplus_{n\geq n_1} \underline{M}'_n\),
